\documentclass[11pt]{article}
\usepackage[margin=1in]{geometry}
\usepackage{amsmath,amssymb,amsthm}
\usepackage{algorithm}
\usepackage{algpseudocode}
\usepackage{booktabs}
\usepackage{hyperref}

\newtheorem{definition}{Definition}
\newtheorem{theorem}{Theorem}
\newtheorem{lemma}[theorem]{Lemma}
\newtheorem{property}{Property}

\title{\textbf{Mathematical Specification, Conflict Orthogonalization, and Homoscedastic Uncertainty Weighting of Multi-Task Losses (ALGO-NN-22)}}
\author{Team Scaibu \\ \small Email: \texttt{jaydeep.9664920749@gmail.com} \\ \small Architecture and Core Engine Team}
\date{\today}

\begin{document}
\maketitle

\begin{abstract}
Joint optimization of heterogeneous tasks often suffers from gradient dominance and destructive interference (negative inner products between task gradients). We analyze two complementary solutions: Bayesian Homoscedastic Uncertainty Weighting, which learns variance-regularized task scales, and Projected Conflicting Gradients (PCGrad), which geometrically orthogonalizes competing update vectors. We derive the projection operators, prove the Non-Destructive Multi-Task Descent Theorem, and provide complete pseudocode matching the reference implementation.
\end{abstract}

\section{Notation and Mathematical Preliminaries}
\label{sec:notation}

Let $K \in \mathbb{N}$ denote the number of simultaneous tasks producing losses $\mathcal{L}_1, \dots, \mathcal{L}_K \in \mathbb{R}_{\ge 0}$. Let $\theta \in \mathbb{R}^P$ denote shared neural network parameters.

\subsection{Task Uncertainty and Gradient Projection Formulations}
\paragraph{Intuitive Meaning:} When training a model for both classification and depth estimation, their loss magnitudes differ by orders of magnitude. Uncertainty weighting learns task variances $\sigma_k^2$ so noisy tasks automatically receive lower weight, while PCGrad projects conflicting gradients to be orthogonal before summing.

\paragraph{Formal Mathematical Definition:}
\begin{itemize}
    \item \textbf{Homoscedastic Uncertainty Weighting:} Introducing learned log-variances $s_k = \log\sigma_k^2$:
    \begin{equation}
    \label{eq:uncertainty_def}
    \mathcal{L}_{\mathrm{total}}(\theta, \mathbf{s}) = \sum_{k=1}^K \left[ \exp(-s_k) \mathcal{L}_k(\theta) + \frac{1}{2} s_k \right].
    \end{equation}
    \item \textbf{PCGrad (Projected Conflicting Gradients):} When task gradients $\mathbf{g}_i, \mathbf{g}_j \in \mathbb{R}^P$ satisfy $\mathbf{g}_i \cdot \mathbf{g}_j < 0$:
    \begin{equation}
    \label{eq:pcgrad_def}
    \mathbf{g}_i \gets \mathbf{g}_i - \frac{\mathbf{g}_i \cdot \mathbf{g}_j}{\|\mathbf{g}_j\|_2^2} \mathbf{g}_j.
    \end{equation}
\end{itemize}

\subsection{Summary of Symbols}
\begin{itemize}
    \item $K \in \mathbb{N}$: Number of concurrent task losses.
    \item $P \in \mathbb{N}$: Number of shared parameter coordinates.
    \item $\mathbf{L} = [L_1, \dots, L_K]^\top \in \mathbb{R}_{\ge 0}^K$: Individual task losses.
    \item $\mathbf{s} = [s_1, \dots, s_K]^\top \in \mathbb{R}^K$: Task log-variance parameters.
    \item $\mathbf{G} \in \mathbb{R}^{K \times P}$: Task gradient vectors.
    \item $\tilde{\mathbf{g}} \in \mathbb{R}^P$: Final aggregated conflict-free gradient vector.
\end{itemize}

\section{Problem Definition and Core Concepts}
\label{sec:problem_definition}

\subsection{Geometric Normal Plane Projection}
The projection operator $\operatorname{Proj}_{\mathbf{n}}(\mathbf{g})$ removes the parallel component along $-\mathbf{n}$, ensuring the projected vector lies strictly in the half-space $\{\mathbf{v} \in \mathbb{R}^P \mid \mathbf{v} \cdot \mathbf{n} \ge 0\}$.

\section{Algorithmic Specification}
\label{sec:algorithm}

Algorithm~\ref{alg:multitask_balancing} specifies the complete execution logic matching \texttt{impl.py}.

\begin{algorithm}[H]
\caption{Multi-Task Loss Balancing and PCGrad Conflict Orthogonalization}
\label{alg:multitask_balancing}
\begin{algorithmic}[1]
\Require Mode $\mathrm{mode} \in \{\text{uncertainty}, \text{pcgrad}\}$, losses $\mathbf{L} \in \mathbb{R}^K$, log-variances $\mathbf{s} \in \mathbb{R}^K$, or gradients $\mathbf{G} \in \mathbb{R}^{K \times P}$.
\Ensure Combined loss $\mathcal{L}_{\mathrm{total}}$, effective weights $\mathbf{w}$, or projected gradient $\tilde{\mathbf{g}} \in \mathbb{R}^P$.
\If{$\mathrm{mode} = \text{uncertainty}$}
    \State $\mathcal{L}_{\mathrm{total}} \gets 0.0, \quad \mathbf{w} \gets []$
    \For{$k = 1$ to $K$}
        \State $w_k \gets \exp(-s_k)$
        \State $\mathcal{L}_{\mathrm{total}} \gets \mathcal{L}_{\mathrm{total}} + (w_k \cdot L_k + 0.5 \cdot s_k)$
        \State Append $w_k$ to $\mathbf{w}$
    \EndFor
    \State \Return $\mathcal{L}_{\mathrm{total}}, \quad \mathbf{w}$
\Else
    \State Copy projected gradient buffer $\tilde{\mathbf{G}} \gets \mathbf{G}$
    \For{$i = 1$ to $K$}
        \For{$j = 1$ to $K$}
            \If{$i \neq j$}
                \State $dot \gets \tilde{\mathbf{G}}_i \cdot \mathbf{G}_j$
                \If{$dot < 0$}
                    \State $norm\_sq \gets \|\mathbf{G}_j\|_2^2 + \epsilon$
                    \State $\tilde{\mathbf{G}}_i \gets \tilde{\mathbf{G}}_i - \frac{dot}{norm\_sq} \mathbf{G}_j$
                \EndIf
            \EndIf
        \EndFor
    \EndFor
    \State $\tilde{\mathbf{g}} \gets \sum_{i=1}^K \tilde{\mathbf{G}}_i$
    \State \Return $\tilde{\mathbf{g}}, \quad \mathbf{w} = [1.0]_{k=1}^K$
\EndIf
\end{algorithmic}
\end{algorithm}

\section{Theoretical Guarantees and Conflict Removal}
\label{sec:guarantees}

\begin{theorem}[Non-Destructive Multi-Task Parameter Updates]
\label{thm:pcgrad_non_destructive}
\textbf{Assumptions:} Let task gradients $\mathbf{g}_1, \dots, \mathbf{g}_K \in \mathbb{R}^P$ be orthogonalized via PCGrad into $\tilde{\mathbf{g}}_1, \dots, \tilde{\mathbf{g}}_K$, with step size $\eta \to 0$.
\textbf{Guarantee:} The aggregate update $\Delta\theta = -\eta \sum_{k=1}^K \tilde{\mathbf{g}}_k$ produces a non-negative directional derivative along every individual task: $\nabla\mathcal{L}_i \cdot \sum_{k=1}^K \tilde{\mathbf{g}}_k \ge 0$ for all $i \in \{1, \dots, K\}$.
\textbf{Proof:}
By construction of the normal projection, for every pair $(i, j)$, $\tilde{\mathbf{g}}_i \cdot \mathbf{g}_j \ge 0$.
Taking the inner product with task gradient $\mathbf{g}_i$:
\begin{equation}
\mathbf{g}_i \cdot \sum_{k=1}^K \tilde{\mathbf{g}}_k = \mathbf{g}_i \cdot \tilde{\mathbf{g}}_i + \sum_{k \neq i} \mathbf{g}_i \cdot \tilde{\mathbf{g}}_k \ge \|\tilde{\mathbf{g}}_i\|_2^2 + 0 \ge 0.
\end{equation}
Thus, no task experiences first-order degradation during gradient descent.
\textbf{Limitations:} PCGrad introduces $\mathcal{O}(K^2 P)$ compute overhead per step.
\end{theorem}

\section{Complexity Analysis}
\label{sec:complexity}

\begin{itemize}
    \item \textbf{Time Complexity:} $\mathcal{O}(K)$ for Uncertainty Weighting; $\mathcal{O}(K^2 \cdot P)$ for PCGrad pairwise projections.
    \item \textbf{Space Complexity:} $\mathcal{O}(K \cdot P)$ for gradient buffers.
\end{itemize}

\section{Worked Numerical Example}
\label{sec:worked_example}

Let $K = 2, P = 2$, $\mathbf{g}_1 = [1.0, 0.0], \mathbf{g}_2 = [-0.5, 1.0]$:
\begin{enumerate}
    \item Dot product: $\mathbf{g}_1 \cdot \mathbf{g}_2 = 1.0(-0.5) + 0.0(1.0) = -0.5 < 0$ (conflict).
    \item $\|\mathbf{g}_2\|_2^2 = (-0.5)^2 + 1.0^2 = 0.25 + 1.0 = 1.25$.
    \item Projected $\tilde{\mathbf{g}}_1$:
    $\tilde{\mathbf{g}}_1 = [1.0, 0.0] - \frac{-0.5}{1.25}[-0.5, 1.0] = [1.0, 0.0] + 0.4[-0.5, 1.0] = [1.0 - 0.2, 0.0 + 0.4] = [0.8, 0.4]$.
    \item Verification: $\tilde{\mathbf{g}}_1 \cdot \mathbf{g}_2 = 0.8(-0.5) + 0.4(1.0) = -0.4 + 0.4 = 0.0$ (strictly orthogonal).
\end{enumerate}

\section{Numerical Considerations and Edge Cases}
\label{sec:numerical}

\begin{itemize}
    \item \textbf{Norm Epsilon Regularization:} $\epsilon = 10^{-12}$ in the projection denominator prevents division by zero when a task gradient is identically zero.
\end{itemize}

\begin{thebibliography}{9}
\bibitem{kendall2018} A.~Kendall, Y.~Gal, and R.~Cipolla, ``Multi-Task Learning Using Uncertainty to Weigh Losses for Scene Geometry and Semantics,'' \emph{CVPR}, 2018.
\bibitem{yu2020} T.~Yu et al., ``Gradient Surgery for Multi-Task Learning,'' \emph{NeurIPS}, 2020.
\end{thebibliography}

\end{document}
